% test-010.tex -- comandos \spnD, \spnM, \spmcd y \spmcm
%
% Compilar:  lualatex-dev test-010.tex
% Validar:   show-pdf-tags --xml test-010.pdf | rnv-wrapp
%            verapdf --flavour ua2 --format text test-010.pdf
%
% Cada caso es un parrafo numerado, para poder referirse a el al escuchar
% el PDF. Los comentarios "doc:" indican lo que documenta el manual; en el
% resto, la lectura se revisa escuchando. Los cuatro trabajan en modo
% matematico y reciben un argumento entre parentesis. Los resultados de
% result=true los calcula el modulo Lua. Las entradas invalidas (que
% detienen la compilacion) no van aqui, entre ellas una lista de un solo
% elemento en \spmcd y \spmcm (se requieren al menos dos). Tampoco el 0
% en \spnD y \spnM, que da siempre error (md-invalid-argument).
\DocumentMetadata
  {
    lang=es-CL,
    pdfversion=2.0,
    pdfstandard=ua-2,
    tagging=on,
    tagging-setup={math/setup={mathml-SE}},
  }
\documentclass{article}
\usepackage[spanish]{babel}
\usepackage{unicode-math}
\usepackage{spintent}
\usepackage{hyperref}
\hypersetup
  {
    pdfdisplaydoctitle = true,
    pdftitle           = {test-010: spnD, spnM, spmcd y spmcm},
  }
\pagestyle{empty}
\begin{document}

\section*{Comando spnD}

% doc: «divisores de n»
1. Sin argumento: $\spnD$.

% doc: «divisores de 6»
2. Con un número: $\spnD(6)$.

% doc: «divisores de m»
3. Con una letra: $\spnD(m)$.

% doc: «divisores de 12»
4. Con el número 12: $\spnD(12)$.

% doc: «divisores de 100»
5. Con el número 100: $\spnD(100)$.

% doc: «divisores de 1»
6. Con el número 1: $\spnD(1)$.

% doc: «divisores de 7»
7. Con el número 7: $\spnD(7)$.

% doc: «divisores de n»
8. Con la letra n: $\spnD(n)$.

% doc: «divisores de a»
9. Con la letra a: $\spnD(a)$.

% doc: «divisores de 6»
10. Con concept true, el valor por defecto: $\spnD[concept=true](6)$.

% doc: «6»
11. Con concept false: $\spnD[concept=false](6)$.

% doc: «m»
12. Con concept false y una letra: $\spnD[concept=false](m)$.

% doc: «los divisores de 6»
13. Con prefix: $\spnD[prefix={los}](6)$.

% doc: «los divisores de m»
14. Con prefix y una letra: $\spnD[prefix={los}](m)$.

% doc: «divisores de 6 son 1 2 3 y 6»
15. Con result true: $\spnD[result=true](6)$.

% doc: «divisores de 12 son 1 2 3 4 6 y 12»
16. Con result true y el 12: $\spnD[result=true](12)$.

% doc: «divisores de 7 son 1 y 7»
17. Con result true y el 7: $\spnD[result=true](7)$.

% doc: «divisores de 1 son 1»
18. Con result true y el 1: $\spnD[result=true](1)$.

% doc: «divisores de 12 son 1 y 2 puntos suspensivos»
19. Con result-num 2 y el 12: $\spnD[result=true,result-num=2](12)$.

% doc: «divisores de 100 son 1 2 4 5 10 20 25 50 y 100»
20. Con result-num 10 y el 100: $\spnD[result=true,result-num=10](100)$.

% doc: «divisores de 12 son 1 puntos suspensivos»
21. Con result-num 1 y el 12: $\spnD[result=true,result-num=1](12)$.

% doc: «divisores de 6»
22. Con result false explícito: $\spnD[result=false](6)$.

% doc: «»
23. Con silent true: $\spnD[silent=true](6)$.

% doc: «divisores de 6»
24. Con silent false explícito: $\spnD[silent=false](6)$.

% doc: «»
25. Con silent true y una letra: $\spnD[silent=true](m)$.

% doc: «multiplicado por»
26. Con silent y result: $\spnD[silent=true,result=true](6)$.

% doc: «6 son 1 2 3 y 6»
27. Con concept false y result: $\spnD[concept=false,result=true](6)$.

% doc: «los divisores de 6 son 1 2 3 y 6»
28. Con prefix y result: $\spnD[prefix={los},result=true](6)$.

% doc: «6»
29. Con prefix y concept false: $\spnD[prefix={los},concept=false](6)$.

% doc: «los divisores de 12 son 1 2 3 puntos suspensivos»
30. Con todas las llaves: $\spnD[prefix={los},concept=true,result=true,result-num=3](12)$.

\section*{Comando spnM}

% doc: «múltiplos de n»
31. Sin argumento: $\spnM$.

% doc: «múltiplos de 6»
32. Con un número: $\spnM(6)$.

% doc: «múltiplos de m»
33. Con una letra: $\spnM(m)$.

% doc: «múltiplos de 12»
34. Con el número 12: $\spnM(12)$.

% doc: «múltiplos de 100»
35. Con el número 100: $\spnM(100)$.

% doc: «múltiplos de 1»
36. Con el número 1: $\spnM(1)$.

% doc: «múltiplos de 7»
37. Con el número 7: $\spnM(7)$.

% doc: «múltiplos de n»
38. Con la letra n: $\spnM(n)$.

% doc: «múltiplos de a»
39. Con la letra a: $\spnM(a)$.

% doc: «múltiplos de 6»
40. Con concept true, el valor por defecto: $\spnM[concept=true](6)$.

% doc: «6»
41. Con concept false: $\spnM[concept=false](6)$.

% doc: «m»
42. Con concept false y una letra: $\spnM[concept=false](m)$.

% doc: «los múltiplos de 6»
43. Con prefix: $\spnM[prefix={los}](6)$.

% doc: «los múltiplos de m»
44. Con prefix y una letra: $\spnM[prefix={los}](m)$.

% doc: «múltiplos de 6 son 6 y 12 puntos suspensivos»
45. Con result true: $\spnM[result=true](6)$.

% doc: «múltiplos de 12 son 12 y 24 puntos suspensivos»
46. Con result true y el 12: $\spnM[result=true](12)$.

% doc: «múltiplos de 1 son 1 y 2 puntos suspensivos»
47. Con result true y el 1: $\spnM[result=true](1)$.

% doc: «múltiplos de 7 son 7 y 14 puntos suspensivos»
48. Con result true y el 7: $\spnM[result=true](7)$.

% doc: «múltiplos de 6 son 6 y 12 puntos suspensivos»
49. Con result-num 2 y el 6: $\spnM[result=true,result-num=2](6)$.

% doc: «múltiplos de 12 son 12 24 36 puntos suspensivos»
50. Con result-num 3 y el 12: $\spnM[result=true,result-num=3](12)$.

% doc: «múltiplos de 6 son 6 12 18 24 30 36 42 48 54 60 puntos suspensivos»
51. Con result-num 10 y el 6: $\spnM[result=true,result-num=10](6)$.

% doc: «múltiplos de 6»
52. Con result false explícito: $\spnM[result=false](6)$.

% doc: «»
53. Con silent true: $\spnM[silent=true](6)$.

% doc: «múltiplos de 6»
54. Con silent false explícito: $\spnM[silent=false](6)$.

% doc: «»
55. Con silent true y una letra: $\spnM[silent=true](m)$.

% doc: «multiplicado por»
56. Con silent y result: $\spnM[silent=true,result=true](6)$.

% doc: «6 son 6 y 12 puntos suspensivos»
57. Con concept false y result: $\spnM[concept=false,result=true](6)$.

% doc: «los múltiplos de 6 son 6 y 12 puntos suspensivos»
58. Con prefix y result: $\spnM[prefix={los},result=true](6)$.

% doc: «6»
59. Con prefix y concept false: $\spnM[prefix={los},concept=false](6)$.

% doc: «los múltiplos de 12 son 12 24 36 puntos suspensivos»
60. Con todas las llaves: $\spnM[prefix={los},concept=true,result=true,result-num=3](12)$.

\section*{Comando spmcd}

% doc: «máximo común divisor»
61. Sin argumento: $\spmcd$.

% doc: «máximo común divisor entre 2 3 y 6»
62. Con tres números: $\spmcd(2, 3, 6)$.

% doc: «máximo común divisor entre m y n»
63. Con dos letras: $\spmcd(m, n)$.

% doc: «máximo común divisor entre 12 y 18»
64. Con los números 12, 18: $\spmcd(12, 18)$.

% doc: «máximo común divisor entre 100 y 75»
65. Con los números 100, 75: $\spmcd(100, 75)$.

% doc: «máximo común divisor entre 7 y 13»
66. Con los números 7, 13: $\spmcd(7, 13)$.

% doc: «máximo común divisor entre 8 12 y 20»
67. Con los números 8, 12, 20: $\spmcd(8, 12, 20)$.

% doc: «máximo común divisor entre 2 3 6 y 12»
68. Con los números 2, 3, 6, 12: $\spmcd(2, 3, 6, 12)$.

% doc: «máximo común divisor entre 48 y 36»
69. Con los números 48, 36: $\spmcd(48, 36)$.

% doc: «máximo común divisor entre 6 y 6»
70. Con los números 6, 6: $\spmcd(6, 6)$.

% doc: «máximo común divisor entre 12 y 18»
71. Sin espacios entre los números: $\spmcd(12,18)$.

% doc: «máximo común divisor entre m n y p»
72. Con las letras m, n, p: $\spmcd(m, n, p)$.

% doc: «máximo común divisor entre a y b»
73. Con las letras a, b: $\spmcd(a, b)$.

% doc: «máximo común divisor entre 2 y m»
74. Con un número y una letra: $\spmcd(2, m)$.

% doc: «12 y 18»
75. Con concept false: $\spmcd[concept=false](12, 18)$.

% doc: «máximo común divisor entre 12 y 18»
76. Con concept true explícito: $\spmcd[concept=true](12, 18)$.

% doc: «máximo común divisor entre 12 y 18»
77. Con read-cmd clear, el valor por defecto: $\spmcd[read-cmd=clear](12, 18)$.

% doc: «mcd entre 12 y 18»
78. Con read-cmd terse: $\spmcd[read-cmd=terse](12, 18)$.

% doc: «mcd entre m y n»
79. Con read-cmd terse y letras: $\spmcd[read-cmd=terse](m, n)$.

% doc: «máximo común divisor entre 12 y 18»
80. Con upper-case false, el valor por defecto: $\spmcd[upper-case=false](12, 18)$.

% doc: «máximo común divisor entre 12 y 18»
81. Con upper-case true: $\spmcd[upper-case=true](12, 18)$.

% doc: «máximo común divisor»
82. Con upper-case true sin argumento: $\spmcd[upper-case=true]$.

% doc: «el máximo común divisor entre 12 y 18»
83. Con prefix: $\spmcd[prefix={el}](12, 18)$.

% doc: «el mcd entre 12 y 18»
84. Con prefix y read-cmd terse: $\spmcd[prefix={el},read-cmd=terse](12, 18)$.

% doc: «máximo común divisor entre 2 3 y 6 es igual a 1»
85. Con result true: $\spmcd[result=true](2, 3, 6)$.

% doc: «máximo común divisor entre 12 y 18 es igual a 6»
86. Con result true y dos números: $\spmcd[result=true](12, 18)$.

% doc: «máximo común divisor entre 2 3 6 y 12 es igual a 1»
87. Con result true y cuatro números: $\spmcd[result=true](2, 3, 6, 12)$.

% doc: «máximo común divisor entre 7 y 13 es igual a 1»
88. Con result true y números primos entre sí: $\spmcd[result=true](7, 13)$.

% doc: «máximo común divisor entre 8 12 y 20 es igual a 4»
89. Con result true y tres números: $\spmcd[result=true](8, 12, 20)$.

% doc: «máximo común divisor entre 2 3 y 6»
90. Con result false explícito: $\spmcd[result=false](2, 3, 6)$.

% doc: «»
91. Con silent true: $\spmcd[silent=true](12, 18)$.

% doc: «máximo común divisor entre 12 y 18»
92. Con silent false explícito: $\spmcd[silent=false](12, 18)$.

% doc: «»
93. Con silent true y letras: $\spmcd[silent=true](m, n)$.

% doc: «el intervalo de a no incluyendo o»
94. Con silent y result: $\spmcd[silent=true,result=true](12, 18)$.

% doc: «12 y 18 es igual a 6»
95. Con concept false y result: $\spmcd[concept=false,result=true](12, 18)$.

% doc: «mcd entre 12 y 18 es igual a 6»
96. Con read-cmd terse y result: $\spmcd[read-cmd=terse,result=true](12, 18)$.

% doc: «máximo común divisor entre 12 y 18 es igual a 6»
97. Con upper-case y result: $\spmcd[upper-case=true,result=true](12, 18)$.

% doc: «el mcd entre 12 y 18 es igual a 6»
98. Con todas las llaves: $\spmcd[prefix={el},read-cmd=terse,upper-case=true,result=true](12, 18)$.

\section*{Comando spmcm}

% doc: «mínimo común múltiplo»
99. Sin argumento: $\spmcm$.

% doc: «mínimo común múltiplo entre 2 3 y 6»
100. Con tres números: $\spmcm(2, 3, 6)$.

% doc: «mínimo común múltiplo entre m y n»
101. Con dos letras: $\spmcm(m, n)$.

% doc: «mínimo común múltiplo entre 12 y 18»
102. Con los números 12, 18: $\spmcm(12, 18)$.

% doc: «mínimo común múltiplo entre 100 y 75»
103. Con los números 100, 75: $\spmcm(100, 75)$.

% doc: «mínimo común múltiplo entre 7 y 13»
104. Con los números 7, 13: $\spmcm(7, 13)$.

% doc: «mínimo común múltiplo entre 8 12 y 20»
105. Con los números 8, 12, 20: $\spmcm(8, 12, 20)$.

% doc: «mínimo común múltiplo entre 2 3 6 y 12»
106. Con los números 2, 3, 6, 12: $\spmcm(2, 3, 6, 12)$.

% doc: «mínimo común múltiplo entre 48 y 36»
107. Con los números 48, 36: $\spmcm(48, 36)$.

% doc: «mínimo común múltiplo entre 6 y 6»
108. Con los números 6, 6: $\spmcm(6, 6)$.

% doc: «mínimo común múltiplo entre 12 y 18»
109. Sin espacios entre los números: $\spmcm(12,18)$.

% doc: «mínimo común múltiplo entre m n y p»
110. Con las letras m, n, p: $\spmcm(m, n, p)$.

% doc: «mínimo común múltiplo entre a y b»
111. Con las letras a, b: $\spmcm(a, b)$.

% doc: «mínimo común múltiplo entre 2 y m»
112. Con un número y una letra: $\spmcm(2, m)$.

% doc: «12 y 18»
113. Con concept false: $\spmcm[concept=false](12, 18)$.

% doc: «mínimo común múltiplo entre 12 y 18»
114. Con concept true explícito: $\spmcm[concept=true](12, 18)$.

% doc: «mínimo común múltiplo entre 12 y 18»
115. Con read-cmd clear, el valor por defecto: $\spmcm[read-cmd=clear](12, 18)$.

% doc: «mcm entre 12 y 18»
116. Con read-cmd terse: $\spmcm[read-cmd=terse](12, 18)$.

% doc: «mcm entre m y n»
117. Con read-cmd terse y letras: $\spmcm[read-cmd=terse](m, n)$.

% doc: «mínimo común múltiplo entre 12 y 18»
118. Con upper-case false, el valor por defecto: $\spmcm[upper-case=false](12, 18)$.

% doc: «mínimo común múltiplo entre 12 y 18»
119. Con upper-case true: $\spmcm[upper-case=true](12, 18)$.

% doc: «mínimo común múltiplo»
120. Con upper-case true sin argumento: $\spmcm[upper-case=true]$.

% doc: «el mínimo común múltiplo entre 12 y 18»
121. Con prefix: $\spmcm[prefix={el}](12, 18)$.

% doc: «el mcm entre 12 y 18»
122. Con prefix y read-cmd terse: $\spmcm[prefix={el},read-cmd=terse](12, 18)$.

% doc: «mínimo común múltiplo entre 2 3 y 6 es igual a 6»
123. Con result true: $\spmcm[result=true](2, 3, 6)$.

% doc: «mínimo común múltiplo entre 12 y 18 es igual a 36»
124. Con result true y dos números: $\spmcm[result=true](12, 18)$.

% doc: «mínimo común múltiplo entre 2 3 6 y 12 es igual a 12»
125. Con result true y cuatro números: $\spmcm[result=true](2, 3, 6, 12)$.

% doc: «mínimo común múltiplo entre 7 y 13 es igual a 91»
126. Con result true y números primos entre sí: $\spmcm[result=true](7, 13)$.

% doc: «mínimo común múltiplo entre 8 12 y 20 es igual a 120»
127. Con result true y tres números: $\spmcm[result=true](8, 12, 20)$.

% doc: «mínimo común múltiplo entre 2 3 y 6»
128. Con result false explícito: $\spmcm[result=false](2, 3, 6)$.

% doc: «»
129. Con silent true: $\spmcm[silent=true](12, 18)$.

% doc: «mínimo común múltiplo entre 12 y 18»
130. Con silent false explícito: $\spmcm[silent=false](12, 18)$.

% doc: «»
131. Con silent true y letras: $\spmcm[silent=true](m, n)$.

% doc: «el intervalo de a no incluyendo o»
132. Con silent y result: $\spmcm[silent=true,result=true](12, 18)$.

% doc: «12 y 18 es igual a 36»
133. Con concept false y result: $\spmcm[concept=false,result=true](12, 18)$.

% doc: «mcm entre 12 y 18 es igual a 36»
134. Con read-cmd terse y result: $\spmcm[read-cmd=terse,result=true](12, 18)$.

% doc: «mínimo común múltiplo entre 12 y 18 es igual a 36»
135. Con upper-case y result: $\spmcm[upper-case=true,result=true](12, 18)$.

% doc: «el mcm entre 12 y 18 es igual a 36»
136. Con todas las llaves: $\spmcm[prefix={el},read-cmd=terse,upper-case=true,result=true](12, 18)$.

\section*{En fórmulas}

% doc: «divisores de 6 es un subconjunto de divisores de de 12»
137. Divisores incluidos: $\spnD(6) \subset \spnD(12)$.

% doc: «divisores de 6 son 1 2 3 y 6 es un subconjunto de divisores de de 12 son 1 2 3 4 6 y 12»
138. Divisores con resultado: $\spnD[result=true](6) \subset \spnD[result=true](12)$.

% doc: «múltiplos de 4 son 4 8 12 puntos suspensivos intersección múltiplos de de 6 son 6 12 18 puntos suspensivos»
139. Múltiplos con resultado: $\spnM[result=true,result-num=3](4) \cap \spnM[result=true,result-num=3](6)$.

% doc: «máximo común divisor entre m y n por mínimo común múltiplo entre m y n es igual a m por n»
140. Relación entre mcd y mcm: $\spmcd(m, n) \spmsym \spmcm(m, n) = m \spmsym n$.

% doc: «máximo común divisor entre 12 y 18 es igual a 6 más mínimo común múltiplo entre 12 y 18 es igual a 36»
141. Con números y resultado: $\spmcd[result=true](12, 18) + \spmcm[result=true](12, 18)$.

% doc: «máximo común divisor entre 12 y 18 por mínimo común múltiplo entre 12 y 18 es igual a 12 por 18»
142. Con otros comandos de la serie: $\spmcd(12, 18) \spmsym \spmcm(12, 18) = \spmul{12}{18}$.

% doc: «divisores de 100 son 1 2 4 5 10 20 25 50 y 100»
143. En fórmula destacada:
\[ \spnD[result=true,result-num=10](100) \].

\end{document}
